% =====================================================================
%  Section 4 -- The moduli space and its hypotheses
%  reader-integrity sentinel — non-ASCII, do not remove: é ö ±
%  CT-S37-DRAFT-B.  Draft only.
%
%  OUTLINE revision 6, repair R-(c): the family (4.1), the flow (4.4)
%  and the Hamiltonian H_V (4.5) are PRINTED here, not cited, because
%  the uniqueness argument of the following section rests on the
%  regularity of (4.4).  Every equation number in parentheses is van
%  der Put--Saito's own.
% =====================================================================

\section{The moduli space and its hypotheses}\label{sec:moduli}

\subsection{Type \texorpdfstring{$(0,0,1,-)$}{(0,0,1,-)} and the strong
Riemann--Hilbert problem}\label{ssec:type001}

For a differential module of type $(0,0,1,-)$ the \emph{strong} Riemann--Hilbert
problem has a positive answer, and the module is represented by
\begin{equation}\label{eq:module}
  \frac{d}{dz}+\frac{A_0}{z}+\frac{A_1}{z-1}+A_\infty ,
\end{equation}
with $A_0,A_1$ constant trace-zero $2\times2$ matrices, $A_\infty$ normalised
to
\begin{equation}\label{eq:Ainf}
  A_\infty=\frac{t}{2}\operatorname{diag}(-1,1),\qquad t\in\mathbb C^{*},
\end{equation}
with $\det A_i=-\theta_i^2/4$, and with the singular points normalised to
$0,1,\infty$. This is a reading of van der Put--Saito \S3.2, p.~2637, from the
page image, and it is \textbf{INSPECTED} (\texttt{S16-01}).

The same reading establishes a negative fact that the paper relies on:
\emph{no semisimplicity hypothesis and no non-resonance hypothesis appears
anywhere in this definition} (\texttt{S16-01}, \textbf{INSPECTED}).

\subsection{Resonance and reducibility are not excluded}\label{ssec:resonance}

p.~2636 defines resonance as a difference of eigenvalues lying in
$\mathbb Z\setminus\{0\}$, and presents resonance and reducibility as sources
of \emph{singularities of the monodromy space} $R$. Modules with such local
data are therefore \emph{inside} the moduli space, not excluded from it. This
too is \textbf{INSPECTED} (\texttt{S16-02}).

The consequence matters for \S\ref{sec:systems}: the logarithmic local data
established in \S\ref{ssec:nonapparent} --- exponents $\{0,0\}$ with nonzero
log coefficient --- are admissible here, although they are excluded by the
non-resonance hypothesis of the more familiar Jimbo normalisation. The
framework used in this paper is van der Put--Saito's, and the difference is
stated rather than smoothed over.

\subsection{The explicit family \texorpdfstring{$(4.1)$}{(4.1)}}
\label{ssec:family}

With $D(z)=z(z-1)$ and
\[
  A=\begin{pmatrix} a & b\\ c & -a\end{pmatrix},
  \qquad y'=Ny,\qquad N=-A/D,
\]
the family $(4.1)$ of p.~2655 is
\begin{align}
  a &= p+\tfrac12\bigl(-2qt+t-\theta_\infty\bigr)(z-q)
        -\tfrac12\,t\,(z-q)^2, \label{eq:fam-a}\\[2pt]
  b &= -\frac{z\Bigl(\bigl(p+\tfrac12(q-1)q\,t\bigr)^2-\dfrac{\theta_1^2}{4}
        +(q-1)\Bigl(\dfrac{\theta_0^2}{4}-\dfrac{\theta_1^2}{4}
        -q\,\dfrac{\theta_\infty^2}{4}\Bigr)\Bigr)}{q} \notag\\
    &\qquad -\frac{(q-1)\bigl(p+\tfrac12((q-1)q\,t-\theta_0+q\theta_\infty)\bigr)
        \bigl(p+\tfrac12((q-1)q\,t+\theta_0+q\theta_\infty)\bigr)}{q},
        \label{eq:fam-b}\\[2pt]
  c &= -\frac{z-q}{q-1}. \label{eq:fam-c}
\end{align}

\begin{remark}[A typo in the printed source]\label{rem:typo}
On p.~2655 the first square in the entry $b(z)$ prints
$\bigl(p+\tfrac12(q-1)q\bigr)^2$, with no $t$, while the second term of the
same entry carries $(q-1)q\,t$. Tested against the source's \emph{own}
constraint $\det A(1)=-\theta_1^2/4$, the printed form leaves a nonzero
remainder, whereas the reading $\bigl(p+\tfrac12(q-1)q\,t\bigr)^2$ used in
\eqref{eq:fam-b} satisfies it identically; the companion constraint
$\det A(0)=-\theta_0^2/4$ holds for both readings and therefore does not
discriminate. The discrimination is by an exact symbolic residual, and the
correction is \textbf{VERIFIED} (\texttt{S16-05}). It is flagged here rather
than adopted silently.

A second discrepancy in the same neighbourhood is recorded and \emph{not}
used: the prose statement that $\theta_\infty$ is the $(1,1)$ entry of
$A_0+A_1$ does not hold for the printed family, where that entry is
$-\theta_\infty/2$. Nothing in this paper depends on it.
\end{remark}

\subsection{The chart \texorpdfstring{$U_0$}{U0}, the Hamiltonian
\texorpdfstring{$(4.5)$}{(4.5)} and the flow
\texorpdfstring{$(4.4)$}{(4.4)}}\label{ssec:flow}

The chart is
\begin{equation}\label{eq:chart}
  U_0\;\cong\;\mathbb C\times(\mathbb C\setminus\{0,1\}),
  \qquad\text{coordinates }(p,q),\ q\notin\{0,1\}.
\end{equation}
On $U_0\times(\mathbb C\setminus\{0\})$ the foliation given by the $2$-form
\begin{equation}\label{eq:omega}
  \Omega=\frac{dp\wedge dq}{q(q-1)}-dH_V\wedge dt
\end{equation}
\emph{coincides with the foliation given by isomonodromy}. That sentence is
quoted from p.~2656 and is \textbf{INSPECTED} (\texttt{S16-03}); the flow and
the Hamiltonian were cropped from the page image and transcribed by looking,
not by optical character recognition. They are
\begin{align}
  (4.4)\qquad
  &\frac{dp}{dt}=q(q-1)\,\frac{\partial H_V}{\partial q},
   \qquad
   \frac{dq}{dt}=-q(q-1)\,\frac{\partial H_V}{\partial p},
   \label{eq:flow}\\[4pt]
  (4.5)\qquad
  &H_V(p,q,t)=-\frac{p^2}{(q-1)q\,t}
     -\frac{\theta_0^2}{4qt}
     +\frac{\theta_1^2}{4(q-1)t}
     +\tfrac14\,q\bigl(qt-t+2\theta_\infty-2\bigr).
   \label{eq:HV}
\end{align}

Both are printed in full, and not merely cited, because the uniqueness argument
of the following section is an argument about \eqref{eq:flow}: if
$q(q-1)\neq0$ and $t\neq0$ at a point of $U_0\times(\mathbb C\setminus\{0\})$,
then the vector field
$v=\partial_t+v_p\partial_p+v_q\partial_q$ determined by \eqref{eq:flow} is
regular there, and exactly one leaf passes through that point. Regularity is
the reason a deformation exists; it is not an ansatz. The verification of
these two non-vanishing conditions at the located point belongs to the
following section and is not performed here.

\subsection{What the time variable is}\label{ssec:time}

The three singular points of \eqref{eq:module} are normalised to $0,1,\infty$,
so $t$ is not the position of a singular point. By \eqref{eq:Ainf}, $t$ is the
\emph{scale of the irregular type at infinity}. This is why an object of type
$(0,0,1,-)$ carries a fixed $t$ with no deformation parameter written alongside
it: $t$ is an invariant computed from the configuration. The reading of the
normalisation is \textbf{INSPECTED} (\texttt{S16-01}); the identification of
what $t$ is, being an interpretation of that normalisation rather than a
quotation, is \textbf{STRUCTURAL} (\texttt{S16-04}).

\subsection{Painlev\'e conventions fixed by citation}\label{ssec:conventions}

Three conventions are fixed by reading printed sources, and each is
\textbf{INSPECTED}; none is re-derived here.

\begin{enumerate}
  \item The Painlev\'e~V equation is taken in the DLMF~32.2.5 convention, with
        parameters $(\alpha,\beta,\gamma,\delta)$
        (\texttt{CT-S9-PVPARAMS-C002}).
  \item The dictionary between those parameters and the monodromy exponents
        $(\theta_0,\theta_t,\theta_\infty)$ is Jimbo's, p.~1152:
        \begin{equation}\label{eq:jimbo}
          \alpha=\tfrac18(\theta_0-\theta_t+\theta_\infty)^2,\quad
          \beta=-\tfrac18(\theta_0-\theta_t-\theta_\infty)^2,\quad
          \gamma=1-\theta_0-\theta_t,\quad
          \delta=-\tfrac12
        \end{equation}
        (\texttt{CT-S9-PVPARAMS-C003}).
  \item Van der Put--Saito's Table~2.1 contains the row labelled
        $\widetilde D_5$/PV, whose Katz invariants are $(0,0,1,-)$
        (\texttt{CT-S9-PVPARAMS-C004}; the folio and the selection rule are
        recorded at \texttt{S6-C1}).
\end{enumerate}

Two warnings attach to \eqref{eq:jimbo}. First, Jimbo's normalisation carries a
non-resonance hypothesis that van der Put--Saito's does not
(\S\ref{ssec:resonance}), so the dictionary is a translation of parameters and
not an importation of hypotheses. Second, by
Remark~\ref{rem:necessary} a multiset match is a necessary condition only, so
\eqref{eq:jimbo} is never used in this paper to \emph{assign} a Painlev\'e
label; it is used to state what a label would have to mean.

\subsection{What membership would require}\label{ssec:membership}

Collecting the above: to place a given operator in this moduli space one must
exhibit a point $(p,q)\in U_0$ and a $t\in\mathbb C^{*}$ at which the family
\eqref{eq:fam-a}--\eqref{eq:fam-c} reproduces that operator's projective normal
form. Such a location is a numerical statement about an explicit algebraic
point; it is not a proof of family membership, and the paper does not report it
as one. The located point and the strength of the evidence for it are the
subject of the following section, where the membership conclusion carries the
label \textbf{STRUCTURAL} (\texttt{S16-14}) and expressly does not rest on the
pre-registered monodromy-and-Stokes falsifier at $30$ digits, which was not
run.
